\section{Validation Audits}
\label{app:validation_audits_submission}

The validation design has to address distinct mechanical
channels. First, frequent participants have more chances to
appear near any defendant, so raw co-bidding rates are not
enough. Second, firms may operate in the same procurement
environments as CADE anchors, giving them more opportunities to
meet those anchors even under ordinary competition. Third,
item-level labels can mechanically reuse the same participation
history that constructs the score. The main text reports the
headline results; this appendix records the audit logic and the
interpretation each audit licenses.

\subsection{Participation-Stratified Permutation}
\label{app:permutation_submission}

The conservative CADE benchmark fixes the direct-defendant list
using cases adjudicated before the end of
$\valConservativeCutoff$. Within the always-loser stratum, we
then compare the observed frequent-loser share among cobidders
with a permutation distribution that preserves participation
volume. Specifically, firms are grouped by participation-volume
quartile defined within the always-loser stratum, cobidder
status is randomly permuted within those quartiles, and the
frequent-loser share among placebo cobidders is recomputed
$\valPermPilot$ times. The number of cobidders per quartile is
preserved across permutations by construction; the within-
quartile shuffle reassigns the cobidder marker over the same
set of firms, leaving the per-quartile marginal unchanged. The
$p$-value reports the share of permutations whose placebo
frequent-loser share is at least as extreme as the observed
share.

The observed share is $\valCADEenrich$, compared with a
permutation mean of $\valCADEbase$. The ratio is
$\valMultThreeTwo\!\times$, with $p<0.001$, and the conservative
ROC AUC is $\valAUCprePost$ (95\% CI $\valAUCprePostCI$,
computed via the DeLong method). The permutation answers the
participation-volume concern: the frequent-loser concentration
near CADE defendants is larger than activity volume alone
predicts. It does not by itself address opportunity-set
exposure, which is the purpose of the next audit.

\subsection{Exposure-Discipline Battery}
\label{app:market_opportunity_submission}

The participation permutation does not by itself address market
opportunity. Firms can appear near CADE anchors because they buy
or sell in the same procurement environments, not because their
loser-side behavior is informative for triage. The submitted
evidence therefore treats exposure discipline as a battery of
audits rather than as a single dispositive adjustment.

The verified audit trail contains four components. First, the
sham-FL permutation randomly reassigns the frequent-loser flag
inside the always-loser pool and recomputes AUC against
cobidder labels across $B=2{,}000$ draws. The script is
\texttt{scripts/25\_sham\_fl\_permutation.R}; the output is
\texttt{output/sham\_fl/sham\_summary.csv}. The unit is the firm;
the retained support is the $\valAlwaysLosers$ always-loser
firms; the result is a sham mean of $\valShamMean$, SD
$\valShamSD$, 99th percentile $\valShamQNinetyNine$, and
permutation $p$ $\valShamP$ against the observed FL14 AUC. This
audit addresses participation-volume opportunity. The price
coefficient under the same sham distribution does not reject the
null---observed $+0.064$ against a sham mean of $+0.144$,
$p=\valShamPriceP$---which is consistent with the scope-not-damages
reading of the price evidence in
Section~\ref{sec:price_scope_conclusion}.

Second, the leakage audit holds cobidder firms out of score
construction within folds and reports the out-of-fold AUC. The
script is \texttt{scripts/40\_leakage\_audit\_d3.R}; the output
is \texttt{output/leakage\_audit\_d3/leakage\_audit\_d3.csv}.
The unit is the firm in the fold construction and the item in
the item-level reference row. This audit addresses label reuse.

Third, the temporal and direct-defendant scope audits change the
information set or the target. The strict timing rows use
pre-window participation to rank later cobidder exposure; the
direct-defendant rows change the positive class to direct CADE
defendants. These audits are reported from
\texttt{output/strict\_train\_threshold/strict\_train\_threshold.csv},
\texttt{output/stratum\_scope/stratum\_scope\_metrics.csv}, and
\texttt{output/auc\_direct\_cade/auc\_direct\_cade.csv}. They
address timing and the risk that the loser-side ranking is a
generic CADE-proximity or cartel-membership screen.

Fourth, the within-FL profile balance audit compares cobidders
with non-cobidder frequent losers on deployment breadth,
repeat-pair structure, direct-CADE-facing share, and item-group
portfolio concentration. The script is
\texttt{scripts/60\_theory\_validation\_bridge.R}; the output is
\texttt{output/theory\_bridge/standardized\_diffs.csv}. The unit
is the firm, and the within-FL comparison uses cobidders
($N=191$ in that output) against non-cobidder frequent losers
($N=2{,}544$). This audit addresses whether the target collapses
to a homogeneous high-volume frequent-loser pool.

The battery does \emph{not} implement a standalone
product-by-buyer-by-year-by-modality matched recomputation of the
cobidder AUC. That stronger opportunity-set design is not part of
the verified active audit trail. The interpretation is therefore
limited: the audits substantially reduce the concern that the
ranking is only participation volume, label reuse, timing, or a
generic direct-defendant detector; they do not prove agreement,
identify cartel members, estimate damages, or eliminate all
market-opportunity confounding.

\begin{table}[htbp]
\centering
\caption{Validation Audits and Threats Addressed}
\label{tab:validation_audit_map_submission}
\footnotesize
\begin{adjustbox}{max width=\textwidth,center}
\begin{threeparttable}
\begin{tabular}{L{3.2cm}L{3.7cm}L{2.4cm}L{3.8cm}L{3.2cm}L{3.7cm}}
\toprule
Audit & Threat addressed & Unit & What is held fixed or withheld & Main result or location & Interpretation \\
\midrule
Participation-stratified permutation &
Participation volume alone creates cobidder concentration &
Firm &
Participation-volume quartile within always-losers &
$\valMultThreeTwo\!\times$ excess; AUC $\valAUCprePost$ &
Concentration exceeds volume-preserving placebo assignment \\
Exposure-discipline battery &
Market opportunity to meet CADE anchors drives adjacency &
Firm and item, depending on audit &
Always-loser support, held-out cobidder firms, pre-window scores, changed positive class, and within-FL profile comparisons &
Sham, leakage, timing, scope, and profile-balance outputs; see Appendix D.2 &
Disciplines, but does not eliminate, opportunity-set exposure \\
Leakage audit &
Score construction reuses firms that define item-level labels &
Item/firm &
Cobidder firms withheld from score construction within folds &
Out-of-fold AUC $\valLeakCVAUC$ &
Signal survives the most direct label-reuse channel \\
Temporal holdout &
Retrospective score uses later information &
Firm/year &
Score formed before the test year &
AUC $\valAUCFLfirmTemp$; Table~\ref{tab:temporal_holdout_year_submission} &
Ranking remains separated from random classification out of time \\
Direct-defendant scope check &
Loser-side ranking is misread as generic cartel-membership detection &
Firm/item &
Target changed from cobidders to direct CADE defendants &
Direct-defendant AUC $\valAUCdirectStd$ &
Near-random direct-defendant classification is expected under loser-side scope \\
\bottomrule
\end{tabular}
\begin{tablenotes}
\footnotesize
\item \textit{Notes:} The audits validate forensic priority, not
cartel membership, liability, damages, or proof. ``Held fixed''
and ``withheld'' describe the audit discipline each exercise
adds relative to the baseline adjudication benchmark.
\end{tablenotes}
\end{threeparttable}
\end{adjustbox}
\end{table}

\subsection{Leakage Decomposition}
\label{app:leakage_submission}

The in-sample item-level AUC is intentionally not treated as the
main validation number. It is $\valLeakRefAUC$, and part of that
performance is mechanical: the same firms can help define both
the participation score and the item-level cobidder label. The
audit decomposes that headline into a structural component and a
leakage component.

\begin{table}[htbp]
\centering
\caption{Leakage Audit}
\label{tab:leakage_audit_submission}
\footnotesize
\begin{adjustbox}{max width=\textwidth,center}
\begin{threeparttable}
\begin{tabular}{L{4.4cm}L{5.4cm}cc}
\toprule
Audit & What changes & AUC & 95\% CI \\
\midrule
In-sample reference & Full-sample score predicts any-cobidder item label
& $\valLeakRefAUC$ & $\valLeakRefAUCCI$ \\
Direct-defendant scope check & Same score predicts direct-CADE item label
& $\valLeakDirectAUC$ & $\valLeakDirectAUCCI$ \\
Cobidder-firm holdout & Held-out cobidders are removed from score construction within folds
& $\valLeakCVAUC$ & $\valLeakCVAUCCI$ \\
Temporal holdout & $2009$--$2016$ score predicts $2017$--$2019$ cobidder labels
& $\valAUCFLfirmTemp$ & $\valAUCFLfirmTempCI$ \\
Temporal direct-defendant check & Same temporal score predicts direct-CADE labels
& $\valLeakDirectTempAUC$ & $\valLeakDirectTempAUCCI$ \\
\bottomrule
\end{tabular}
\begin{tablenotes}
\footnotesize
\item \textit{Notes:} The reference AUC is partly tautological.
The firm-holdout and temporal rows isolate the component that
survives when labels are withheld from score construction or
placed out of time. Direct-defendant rows check whether the
loser-side statistic accidentally behaves like a general
membership detector.
Direct-defendant AUCs differ slightly across firm-level,
item-level, and temporal audit definitions, but remain near
random in each specification.
\end{tablenotes}
\end{threeparttable}
\end{adjustbox}
\end{table}

The surviving component is large: AUC remains in the
$\valLeakStructLow$--$\valLeakStructHigh$ range under the
holdout audits. The lost component, roughly
$\valLeakTautLow$--$\valLeakTautHigh$ AUC, is disclosed as
mechanical. That disclosure is important because the main claim
is an evidence-allocation claim. A cheap record-layer statistic
retains meaningful ranking power after the most obvious
mechanical channels are removed.

\subsection{Direct-Defendant Scope Check}
\label{app:direct_scope_submission}

The direct-defendant result is a falsification of overreach. If
the statistic were a generic cartel-membership detector, it
should classify direct CADE defendants as well as cobidders. It
does not. The direct-defendant AUC is $\valAUCdirectStd$, close
to random classification. The item-level direct-defendant rows
in Table~\ref{tab:leakage_audit_submission} use different
samples and labels, but they are likewise near random. That
pattern is informative: the
statistic behaves like a loser-side ranking, while direct
defendants are legal anchors that often appear on the winner side
of the arrangement.

A complementary universe-anchored scope matrix systematically
varies the universe (always-loser firms vs full BEC firm registry)
and the positive class (cobidders vs direct defendants) to map
where the ranking targets succeed and where they fail. The matrix
produces a stronger result than a stand-alone null: raw
participation count on the full BEC universe against direct CADE
defendants returns AUC $\valAUCParticDefendants$, which is
\emph{below} random. A high-participation BEC firm is less likely,
not more likely, to be a direct CADE defendant under the
loser-side scoring logic. The score does not simply fail to rank
winner-heavy defendants; the relevant direction is reversed,
consistent with the role separation observed in D4 (14.9\% of
direct defendants are always-losers; median win rate $0.261$
versus $0.086$ for cobidders).

\begin{table}[htbp]
\centering
\caption{Universe-Anchored Stratum Scope Matrix}
\label{tab:scope_matrix_submission}
\footnotesize
\begin{adjustbox}{max width=\textwidth,center}
\begin{threeparttable}
\begin{tabular}{p{3.3cm}p{3.0cm}p{3.0cm}cp{4.0cm}}
\toprule
Score & Universe & Positive class & AUC & Reading \\
\midrule
Binary FL & Always-losers & Cobidders & $\valScopeRowOne$
  & Within-stratum prioritization \\
Continuous log\_tc & Always-losers & Cobidders & $\valScopeRowTwo$
  & Best within-stratum ranking \\
Binary FL & All BEC firms & Direct CADE & $\valScopeRowThree$
  & At chance; no membership detection \\
Participation count & All BEC firms & Direct CADE & $\valScopeRowFour$
  & \textbf{Below chance: actively repels winner-heavy} \\
Binary FL (frozen 09--16) & Always-losers as of 09--16 & Cobidders
  & $\valScopeRowFive$ & Weaker than pooled within-stratum \\
Continuous (train 09--16) & Items 17--19 & Cobidder items
  & $\valScopeRowSix$ & Some prospective generalization \\
Binary FL (frozen 09--16) & Items 17--19 & Cobidder items
  & $\valScopeRowSeven$ & Weak binary item-level holdout \\
Continuous (train 09--16) & Items 17--19 & Direct-defendant items
  & $\valScopeRowEight$ & Direct-defendant null survives timing \\
\bottomrule
\end{tabular}
\begin{tablenotes}
\footnotesize
\item \textit{Notes:} AUC across an eight-row matrix that varies the
score, the universe of firms or items being ranked, and the positive
class. Row~4 is below chance because the loser-side score
discriminates \emph{against} high-participation winner-heavy
defendants. The matrix structurally rules out a generic-detector
reading.
\end{tablenotes}
\end{threeparttable}
\end{adjustbox}
\end{table}

\subsection{Year-by-Year Holdout}
\label{app:year_holdout_submission}

The rolling-origin exercise trains the score on years preceding
each test year and evaluates cobidder discrimination in the test
year. Table~\ref{tab:temporal_holdout_year_submission} reports the
year-by-year decomposition behind the temporal-holdout figure in
the main text.

\begin{table}[htbp]
\centering
\caption{Temporal-Holdout AUC by Test Year}
\label{tab:temporal_holdout_year_submission}
\footnotesize
\begin{threeparttable}
\begin{tabular}{ccccc}
\toprule
Test year & AUC & 95\% CI & $N$ firms & $N$ CADE \\
\midrule
$2014$ & $0.819$ & $[0.779,\ 0.860]$ & $8{,}631$ & $88$ \\
$2015$ & $0.817$ & $[0.776,\ 0.858]$ & $10{,}122$ & $102$ \\
$2016$ & $0.851$ & $[0.821,\ 0.881]$ & $11{,}699$ & $127$ \\
$2017$ & $0.862$ & $[0.832,\ 0.891]$ & $13{,}057$ & $148$ \\
$2018$ & $0.897$ & $[0.877,\ 0.918]$ & $14{,}391$ & $161$ \\
$2019$ & $0.922$ & $[0.912,\ 0.933]$ & $15{,}597$ & $181$ \\
\bottomrule
\end{tabular}
\begin{tablenotes}
\footnotesize
\item \textit{Notes:} Each row computes the screening score using
participation observed before the test year and evaluates AUC
against the cobidder label in that test year. The exercise is
out-of-time discrimination, not a field deployment.
\end{tablenotes}
\end{threeparttable}
\end{table}

The year-by-year pattern is not a claim that the environment
becomes progressively easier to police. The samples and
adjudication links change across years. The relevant point is
that the score remains separated from random classification in
each rolling-origin test year.

\subsection{Strict-Timing Classifier Battery}
\label{app:timing_battery_submission}

The main text reports the headline strict-timing AUCs. This
subsection records the full three-classifier battery behind them.
Two structural facts make the timing audit harder to dismiss as a
same-firm artifact. First, firm persistence between the
$\valYearStart$--$2016$ and $2017$--$\valYearEnd$ windows is
$\valFirmPersist$: roughly nine in ten always-loser firms in the
test window are firms not seen in the training window, so the
temporal holdout evaluates a substantially fresh firm population
rather than the same firms in different years. Second, when the
training window is shortened further (scores frozen on
$2009$--$2015$, evaluated against cobidders linked to CADE
adjudications closed after $2019$---a strictly out-of-time
target), the firm-level AUC is $\valClfFifteenFL$ for the
frequent-loser flag and $\valClfFifteenPostTC$ for the continuous
score. The ranking remains informative when the classifier is
trained at one horizon, the cobidder label is defined by
adjudications at a later horizon, and the bulk of test-period
firms were not in the training set.

\begin{table}[htbp]
\centering
\caption{Three-Classifier Timing Battery}
\label{tab:three_classifier_submission}
\footnotesize
\begin{adjustbox}{max width=\textwidth,center}
\begin{threeparttable}
\begin{tabular}{p{4.0cm}cccc}
\toprule
& \multicolumn{2}{c}{vs cobid\_all} & \multicolumn{2}{c}{vs cobid\_post2019} \\
\cmidrule(lr){2-3}\cmidrule(lr){4-5}
Classifier & FL flag & Continuous & FL flag & Continuous \\
\midrule
clf\_2015 (train 09--15) & $\valClfFifteenFL$ & $\valClfFifteenTC$
  & $\valClfFifteenPostFL$ & $\valClfFifteenPostTC$ \\
clf\_2017 (train 09--17) & $\valClfSeventeenFL$ & $\valClfSeventeenTC$
  & $\valClfSeventeenPostFL$ & $\valClfSeventeenPostTC$ \\
clf\_2019\_full (in-sample) & $\valAUCFLfirm$ & $\valAUClogtc$
  & --- & --- \\
\bottomrule
\end{tabular}
\begin{tablenotes}
\footnotesize
\item \textit{Notes:} AUC against two cobidder targets across three
classifiers. cobid\_all uses cobidders linked to any CADE adjudication
in the panel. cobid\_post2019 uses cobidders linked only to
adjudications closed after $2019$, a strictly out-of-time target
that cannot be in the clf\_2015 or clf\_2017 training data.
\end{tablenotes}
\end{threeparttable}
\end{adjustbox}
\end{table}
